\BourbakiExerciseGroup

\begin{enumerate}
\def\labelenumi{\arabic{enumi})}
\item
  Find all possible topologies on a set of two elements or three elements.
\item
  \begin{enumerate}
  \def\labelenumii{\alph{enumii})}
  \tightlist
  \item
    Let X be an ordered set. Show that the set of intervals
  \end{enumerate}
\end{enumerate}

\[
[ x, \rightarrow [ \quad (\text { resp. } ] \leftarrow , x ])
\]

is a base of a topology on X; this topology is called the right (resp. left) topology of X. In the right topology, any intersection of open sets is an open set, and the closure of \(\{x\}\) is the interval {]}\(\leftarrow\) , x{]}.

\begin{enumerate}
\def\labelenumi{\alph{enumi})}
\setcounter{enumi}{1}
\item
  A topological space is said to be a Kolmogoroff space if it satisfies the following condition: given any two distinct points \(x\) , \(x'\) of \(X\) , there is a neighbourhood of one of these points which does not contain the other. Show that an ordered set with the right topology is a Kolmogoroff space.
\item
  Let X be a Kolmogoroff space in which every intersection of open sets is an open set. Show that \(x \in \overline{\{x'\}}\) is an order relation between x and \(x'\) in X, and that if this relation is written as \(x \leqslant x'\) then the given topology on X is identical with the right topology determined by this ordering.
\item
  Deduce from c) that if X is a Kolmogoroff space then every finite non-empty subset of X has at least one isolated point. Hence, if X has no isolated point, every non-empty open subset of X is infinite.
\end{enumerate}

\begin{enumerate}
\def\labelenumi{\arabic{enumi})}
\setcounter{enumi}{2}
\tightlist
\item
  If \(\mathbf{A}\) is any subset of a topological space \(\mathbf{X}\) , let \(\alpha(\mathbf{A})\) denote \(\overline{\mathbf{A}}\) and \(\beta(\mathbf{A})\) denote \(\overline{\mathbf{A}}\) . Then \(\mathbf{A} \subset \mathbf{B}\) implies \(\alpha(\mathbf{A}) \subset \alpha(\mathbf{B})\) and \(\beta(\mathbf{A}) \subset \beta(\mathbf{B})\) . \(a)\) Show that if \(\mathbf{A}\) is open then \(\mathbf{A} \subset \alpha(\mathbf{A})\) , and that if \(\mathbf{A}\) is closed then \(\beta(\mathbf{A}) \subset \mathbf{A}\) .
\end{enumerate}

\begin{enumerate}
\def\labelenumi{\alph{enumi})}
\setcounter{enumi}{1}
\item
  Deduce from a) that if A is any subset of X, then \(\alpha(\alpha(A)) = \alpha(A)\) and \(\beta(\beta(A)) = \beta(A)\) .
\item
  Give an example of a set A \(*\) (on the real line) \(_{*}\) such that the seven sets A, \(\mathring{A}\) , \(\overline{A}\) , \(\alpha(A)\) , \(\beta(A)\) , \(\beta(\overline{A})\) , \(\alpha(\mathring{A})\) are all distinct and satisfy no relations of inclusion except the following: \(\mathring{A} \subset A \subset \overline{A}\) , \(\mathring{A} \subset \alpha(\mathring{A}) \subset \beta(A) \subset \overline{A}\) , \(\mathring{A} \subset \alpha(A) \subset \beta(\overline{A}) \subset \overline{A}\) , \(\alpha(\mathring{A}) \subset \alpha(A)\) , \(\beta(A) \subset \beta(\overline{A})\) .
\item
  Show that if U, V are two open sets such that \(U \cap V = \varnothing\) , then also \(\alpha(U) \cap \alpha(V) = \varnothing\) {[}use b){]}.
\end{enumerate}

\begin{enumerate}
\def\labelenumi{\arabic{enumi})}
\setcounter{enumi}{3}
\tightlist
\item
  \begin{enumerate}
  \def\labelenumii{\alph{enumii})}
  \tightlist
  \item
    Give an example * on the real line * of two open sets A, B such that the four sets A ∩ B, B ∩ A̅, A ∩ B and A̅ ∩ B are all distinct.
  \end{enumerate}
\end{enumerate}

\begin{enumerate}
\def\labelenumi{\alph{enumi})}
\setcounter{enumi}{1}
\tightlist
\item
  * Give an example of two intervals A, B on the real line, such that A ∩ B is not contained in A ∩ B. *
\end{enumerate}

\begin{enumerate}
\def\labelenumi{\arabic{enumi})}
\setcounter{enumi}{4}
\tightlist
\item
  If A is any subset of a topological space X, the frontier of A is denoted by \(\operatorname{Fr}(\mathbf{A})\) .
\end{enumerate}

\begin{enumerate}
\def\labelenumi{\alph{enumi})}
\item
  Show that \(\operatorname{Fr}(\overline{\mathbf{A}}) \subset \operatorname{Fr}(\mathbf{A})\) , \(\operatorname{Fr}(\mathring{\mathbf{A}}) \subset \operatorname{Fr}(\mathbf{A})\) , and give an example * on the real line * where these three sets are distinct.
\item
  Let A, B be two subsets of X; show that
\end{enumerate}

\[
\operatorname{Fr} (\mathbf {A} \cup \mathbf {B}) \subset \operatorname{Fr} (\mathbf {A}) \cup \operatorname{Fr} (\mathbf {B})
\]

and give an example * on the real line * where these sets are distinct. If \(\overline{\mathbf{A}} \cap \overline{\mathbf{B}} = \emptyset\) , show that \(\operatorname{Fr}(\mathbf{A} \cup \mathbf{B}) = \operatorname{Fr}(\mathbf{A}) \cup \operatorname{Fr}(\mathbf{B})\) .

\begin{enumerate}
\def\labelenumi{\alph{enumi})}
\setcounter{enumi}{2}
\tightlist
\item
  If \(\mathbf{A}\) and \(\mathbf{B}\) are open in \(\mathbf{X}\) , show that
\end{enumerate}

\[
\begin{array}{r l} & (\mathbf {A} \cap \operatorname{Fr} (\mathbf {B})) \cup (\mathbf {B} \cap \operatorname{Fr} (\mathbf {A})) \subset \operatorname{Fr} (\mathbf {A} \cap \mathbf {B}) \\ & \qquad \subset (\mathbf {A} \cap \operatorname{Fr} (\mathbf {B})) \cup (\mathbf {B} \cap \operatorname{Fr} (\mathbf {A})) \cup (\operatorname{Fr} (\mathbf {A}) \cap \operatorname{Fr} (\mathbf {B})) \end{array}
\]

and give an example * on the real line * where these three sets are distinct.

\begin{enumerate}
\def\labelenumi{\arabic{enumi})}
\setcounter{enumi}{5}
\item
  Show that a subset A of a topological space X meets each dense subset of X if and only if the interior of A is not empty.
\item
  Consider the following four properties of a given topological space \(\mathbf{X}\) :
\end{enumerate}

(D \(_{I}\) ) The topology of X has a countable base.

\((D_{II})\) X has a countable dense subset.

\((D_{III})\) Every subset of X, all of whose points are isolated, is countable.

\((D_{IV})\) Every set of mutually disjoint non-empty open subsets of X is countable.

Show that \((\mathbf{D}_{\mathbf{I}})\) implies \((\mathbf{D}_{\mathbf{II}})\) and \((\mathbf{D}_{\mathbf{III}})\) , and that each of \((\mathbf{D}_{\mathbf{II}})\) and \((\mathbf{D}_{\mathbf{III}})\) imply \((\mathbf{D}_{\mathbf{IV}})\) (*).

\begin{enumerate}
\def\labelenumi{\arabic{enumi})}
\setcounter{enumi}{7}
\item
  If a subset A of a topological space X has no isolated points, then the same is true of its closure \(\overline{A}\) in X.
\item
  Let X be a set and let \(M \to \overline{M}\) be a mapping of \(\mathfrak{P}(X)\) onto itself such that: 1) \(\overline{\varnothing} = \varnothing\) ; 2) for every \(M \subset X\) we have \(M \subset \overline{M}\) ; 3) for every \(M \subset X\) we have \(\overline{M} = \overline{M}\) ; 4) for all \(M \subset X\) and \(N \subset X\) we have
\end{enumerate}

\[
\overline {{\mathbf {M} \cup \mathbf {N}}} = \overline {{\mathbf {M}}} \cup \overline {{\mathbf {N}}}.
\]

Show that there is a unique topology on X such that \(\overline{M}\) is the closure of M with respect to this topology, for all \(M \subset X\) (define the topology by means of its closed sets).

